\documentclass[aspectratio=169,10pt,english]{beamer}

\input{../../../_preambulo_beamer.tex}
\input{../../../_comandos_beamer.tex}

\selectlanguage{english}

\title{MA1001B - Statistical Modeling for Decision Making}
\subtitle{Lesson 18: Geometric and Negative-Binomial Waiting Times}
\author{}
\institute{Tecnol\'ogico de Monterrey}
\date{}

\begin{document}

\begin{frame}[plain]
  \vspace{1.2cm}
  {\Large\bfseries MA1001B - Statistical Modeling for Decision Making\par}
  \vspace{0.7cm}
  {\large Lesson 18: Geometric and Negative-Binomial Waiting Times\par}
  \vspace{0.7cm}
  {\color{TecRojo}\rule{\textwidth}{1.2pt}}
  \vspace{0.8cm}

  MA1001B Faculty\\[0.3cm]
  Term\\[0.3cm]
  Tecnol\'ogico de Monterrey
\end{frame}

\begin{frame}{Waiting for the First Success}
  \begin{alertblock}{Guiding question}
    If each independent audit trial is defective with chance $p=0.12$, how
    long do we wait for the first success, and what happens if fatigue lowers
    $p$?
  \end{alertblock}

  \vspace{0.6em}
  By the end of this lesson, you can:
  \begin{itemize}
    \item compute $P(X=1)=p$ and $E[X]=1/p$;
    \item extend the wait to $r=3$ successes with $E[X]=r/p$;
    \item convert scipy's failure count into a trial count;
    \item explain why declining $p$ makes the mean wait exceed $1/p$.
  \end{itemize}

  \vfill
  \scriptsize All waiting times used today are synthetic. No real company data are used.
\end{frame}

\begin{frame}{Geometric Trials Until the First Defective}
  \small
  $X$ is the trial number of the first success, $X=1,2,3,\ldots$
  \[
    P(X=1)=p=0.12,\qquad E[X]=\frac{1}{p}=8.333333.
  \]
  Verified $P(X\le 5)=0.472268$. scipy \texttt{geom} uses this trials-until-first-success convention.
\end{frame}

\begin{frame}[fragile]{Step 1: Geometric Waiting Time}
  \begin{columns}[T]
    \begin{column}{0.42\textwidth}
      \small
      \begin{lstlisting}
p = 0.12
p1 = stats.geom.pmf(1, p)
mean = 1 / p
      \end{lstlisting}

      \begin{block}{Verified output}
        $P(X=1)=0.120000$\\
        $P(X\le 5)=0.472268$\\
        $E[X]=8.333333$
      \end{block}
    \end{column}
    \begin{column}{0.58\textwidth}
      \centering
      \includegraphics[width=\textwidth]{../figures/waiting_times_01_geometric.png}
      \vspace{0.2em}
      \scriptsize Geometric PMF from Step 1
    \end{column}
  \end{columns}
\end{frame}

\begin{frame}{Negative Binomial: Wait for $r=3$}
  \small
  Now $X$ is the trial of the third success.
  \[
    E[X]=\frac{r}{p}=\frac{3}{0.12}=25,\qquad P(X=3)=p^3=0.001728.
  \]
  scipy \texttt{nbinom} counts failures before $r$ successes. The scripts add
  $r$ so that $X$ is a trial count, matching OpenStax.
\end{frame}

\begin{frame}[fragile]{Step 2: Trials Until $r=3$}
  \begin{columns}[T]
    \begin{column}{0.40\textwidth}
      \small
      \begin{lstlisting}
r, p = 3, 0.12
fail = x - r
mass = stats.nbinom.pmf(
    fail, n=r, p=p
)
      \end{lstlisting}

      \begin{block}{Verified output}
        $E[X]=25.000000$\\
        $\Var(X)=183.333333$\\
        $P(X=3)=0.001728$
      \end{block}
    \end{column}
    \begin{column}{0.60\textwidth}
      \centering
      \includegraphics[width=\textwidth]{../figures/waiting_times_02_negative_binomial.png}
      \vspace{0.2em}
      \scriptsize Negative-binomial PMF from Step 2
    \end{column}
  \end{columns}
\end{frame}

\begin{frame}{Step 3: Fatigue Changes $p$}
  \begin{columns}[T]
    \begin{column}{0.42\textwidth}
      \small
      Let $p_t=0.12\times 0.97^{t-1}$. Then $p_1=0.120000$ but $p_{10}=0.091228$.

      \vspace{0.4em}
      Seed-$42$ mean wait: $17.162800$, versus geometric $8.333333$.

      \vspace{0.4em}
      \textbf{Limit:} trials are not independent if fatigue changes $p$.
    \end{column}
    \begin{column}{0.58\textwidth}
      \centering
      \includegraphics[width=\textwidth]{../figures/waiting_times_03_fatigue_dependence.png}
      \vspace{0.2em}
      \scriptsize Fatigued mean wait from Step 3
    \end{column}
  \end{columns}
\end{frame}

\begin{frame}{Concept Check}
  \begin{enumerate}
    \item Why is $P(X=1)$ equal to $p$, not $(1-p)p$?
    \item Compute $E[X]$ for the geometric model with $p=0.12$.
    \item What is $E[X]$ if we wait for $r=3$ successes?
    \item Why does a declining $p$ make the wait longer than $1/p$?
  \end{enumerate}

  \vspace{0.6em}
  \begin{block}{Connection to Lesson 19}
    The session can continue into continuous random variables, where
    probabilities are areas rather than point masses.
  \end{block}

  \vfill
  \scriptsize Source: OpenStax, \textit{Introductory Business Statistics 2e},
  Section 4.3. CC BY-NC-SA.
\end{frame}

\appendix



\end{document}
